\documentclass[../../../main.tex]{subfiles}
\begin{document}
\subsection{Spin of Single Particle}
\subsubsection{Spin operator.}
There is a fundamental distinction between particles with half-odd integer spin, called Fermions, which obey the Pauli exclusion principle, and those with integer spin, called Bosons, which do not obey the Pauli exclusion principle.
For special case $l=1/2\equiv s$, we introduce operator $\mathbf{S}$, which satisfies
\begin{equation*}
    [\hat{S}_{i}, \hat{S}_{j}] = i \epsilon_{ijk} \hat{S}_{k}
\end{equation*}
The weight of a state $|\chi\rangle$ measures the eigenvalue of the spin operator of the state of the quantum particle
\begin{equation*}
    S_3 |\chi\rangle = m_S |\chi\rangle = \pm \frac{1}{2} |\chi\rangle \quad
    S^2 |\chi\rangle = s(s + 1) |\chi\rangle = \frac{3}{4} |\chi\rangle
\end{equation*}
with
\begin{equation*}
    -s \leq m_S \leq s, \quad  m_S = \pm \frac{1}{2}
\end{equation*}
where $s$ denote the spin of quantum particle.

\subsubsection{spinor.}
The general element of the space spawned by this basis is called a spinor and can be formally written as
\begin{equation*}
    \ket{\boldsymbol{\chi}}= \begin{bmatrix}
        \ket{\chi_+} \\
        \ket{\chi_-}
    \end{bmatrix}
\end{equation*}
In this coordinate system, we have two possible state: one for $m_s=1/2$ and other for $m_s=-1/2$
\begin{equation*}
    \big| \frac{1 }{2 },\frac{1 }{2 }\big\rangle = \ket{\uparrow} = \begin{bmatrix} 1\\0 \end{bmatrix}\qquad
    \big| \frac{1 }{2 },-\frac{1 }{2 }\big\rangle = \ket{\downarrow} = \begin{bmatrix} 0\\1 \end{bmatrix}
\end{equation*}

\subsubsection{Pauli matrices.}
The components of spin $S_i$ can be written in terms of Pauli spin matrix
\begin{equation*}
    \mathbf{S} = \frac{1 }{2 } \boldsymbol{\sigma}
\end{equation*}
with
\begin{equation*}
    \sigma_1 =
    \begin{bmatrix}
        0 & 1 \\
        1 & 0
    \end{bmatrix}
    \quad
    \sigma_2 =
    \begin{bmatrix}
        0 & -i \\
        i & 0
    \end{bmatrix}
    \quad
    \sigma_3    =
    \begin{bmatrix}
        1 & 0  \\
        0 & -1
    \end{bmatrix}
\end{equation*}
such that
\begin{equation*}
    S_1 = \frac{1}{2}
    \begin{bmatrix}
        0 & 1 \\
        1 & 0
    \end{bmatrix}
    \quad
    S_2 = \frac{1}{2}
    \begin{bmatrix}
        0 & -i \\
        i & 0
    \end{bmatrix}
    \quad
    S_3 = \frac{1}{2}
    \begin{bmatrix}
        1 & 0  \\
        0 & -1
    \end{bmatrix}
\end{equation*}
The Pauli matrices satisfy
\begin{gather*}
    \sigma_i^2=\mathbf{1}  \\
    [\sigma_i, \sigma_j] = 2i \epsilon_{ijk} \sigma_k\qquad    \left\{ \sigma_i,\sigma_j \right\} = 2\delta_{ij }\mathbf{1}\\
    (\mathbf{a} \cdot \boldsymbol{\sigma})(\mathbf{b} \cdot \boldsymbol{\sigma}) = \mathbf{a} \cdot \mathbf{b} + i(\mathbf{a} \times \mathbf{b}) \cdot \boldsymbol{\sigma}
\end{gather*}
with vector $\mathbf{a}$ and $\mathbf{b}$.
In components form, the third identity reduce into product identity for Pauli matrices
\begin{equation*}
    \sigma_i \sigma_j = \delta_{ij} \mathbf{1} + i \epsilon_{ijk} \sigma_k
\end{equation*}
Using the Pauli matrices, we can derive the relation for the spin-$1/2$ operator $\mathbf{S} = \boldsymbol{\sigma}/2$ by multiplying previous relation by $1/4$
\begin{equation*}
    S_iS_j= \frac{1 }{4 }\delta_{ij}\mathbf{1} +\frac{i }{2} \epsilon_{ijk }S_k
\end{equation*}
If $i\neq j$, then
\begin{equation*}
    S_iS_j = -S_jS_i
\end{equation*}
Since $S_iS_j = i\epsilon_{ijk }S_k/2$ and $S_jS_i= -i \epsilon_{ijk }S_k/2$.
Along with $\mathbf{1}$, the Pauli matrices span the $2 \times  2$ Hermitian matrices.

\subsubsection{Bloch vector.}
The Bloch vector associated with spinor state $\ket{\chi}$ is
\begin{equation*}
    \mathbf{e }= \braket{\boldsymbol{\sigma}  }_{\chi}= \braket{\chi  | \boldsymbol{\sigma }|\chi}
\end{equation*}
Since the state $\ket{\chi}$ is a pure state, the corresponding Bloch vector represented as a point on the Bloch sphere, with $\mathbf{e}$ being the vector pointing to that location.
Bloch vector can be written in Cartesian components as
\begin{equation*}
    \mathbf{e} = (e_x,e_y,e_z) = (\sin \theta \cos \phi, \sin \theta \sin \phi, \cos \theta)
\end{equation*}
With the spinor parameterization as
\begin{equation*}
    |\chi\rangle = \begin{bmatrix}
        \cos(\theta/2) \\
        e^{i\phi} \sin(\theta/2)
    \end{bmatrix},
\end{equation*}
we can work backward by determining the Bloch vector first the evaluating the spinor, where
\begin{equation*}
    \theta = \arccos(e_z), \quad \varphi = \arctan2(e_y, e_x)
\end{equation*}

\subsubsection{Ladder operator.}
The Ladder operator for spin-$1/2$ particle is written as
\begin{equation*}
    S_\pm = S_1 \pm iS_2
\end{equation*}
or in terms of Pauli matrices
\begin{equation*}
    \sigma_ \pm = \frac{1 }{2 }\left( \sigma_1 \pm \sigma_2 \right)
\end{equation*}
In the dimensionless representation, this mean that $S_\pm=\sigma_\pm$ and their action are
\begin{gather*}
    \sigma_+ |\downarrow\rangle = |\uparrow\rangle, \quad \sigma_- |\uparrow\rangle = |\downarrow\rangle\\
    S_+ |\downarrow\rangle = \frac{1 }{2 } |\uparrow\rangle, \quad S_- |\uparrow\rangle = \frac{1 }{2} |\downarrow\rangle
\end{gather*}
with the extremal states annihilated.
The matrix representation of them are
\begin{equation*}
    \sigma_+ =
    \begin{bmatrix}
        0 & 1 \\
        0 & 0
    \end{bmatrix}
    \quad
    \sigma_- =
    \begin{bmatrix}
        0 & 0 \\
        1 & 0
    \end{bmatrix}
\end{equation*}
They obey the same commutator relation as the angular momentum version, beside the normalization constant
\begin{equation*}
    [\sigma_3, \sigma_\pm] = \pm 2 \sigma_\pm, \quad [\sigma_+, \sigma_-] = \sigma_3
\end{equation*}
and for anticommutator
\begin{equation*}
    \left\{ \sigma_3,\sigma_\pm  \right\} = 0 \quad
    \left\{\sigma_+,\sigma_-  \right\} =\mathbf{1 }
\end{equation*}
The product of ladder operators is a projection operator
\begin{equation*}
    \sigma_+ \sigma_- = \ket{\uparrow}\bra{\uparrow}\qquad
    \sigma_- \sigma_+ = \ket{\downarrow}\bra{\downarrow}
\end{equation*}
This follows naturally by writting the operators in outer product form
\begin{equation*}
    \sigma_+ = \ket{\uparrow }\bra{\downarrow} \qquad \sigma_- = \ket{\downarrow}\bra{\uparrow}
\end{equation*}
In this form, it can be easily seen that they annihilate the extremal state.
Projection operator has idempotent property $P^2=P$ and obey
\begin{equation*}
    f(P) = f(0)I + [f(1) - f(0)]P
\end{equation*}
As such, the ladder operator also obey
\begin{gather*}
    f(\sigma_+\sigma_-) = f(0) + [f(1) - f(0)]\sigma_+\sigma_-    \\
    f(\sigma_-\sigma_+) = f(0) + [f(1) - f(0)]\sigma_+\sigma_+
\end{gather*}
These operators also have similarity transformation
\begin{equation*}
    \exp \left( \frac{\alpha \sigma_z}{2} \right)  \sigma_\pm \exp \left(- \frac{\alpha \sigma_z}{2} \right)  = e^{\pm\alpha} \sigma_\pm
\end{equation*}
with
\begin{equation*}
    \exp \left(\pm \frac{\alpha \sigma_z}{2}  \right)
    = \begin{bmatrix}
        e^{\pm\alpha/2} & 0               \\
        0               & e^{\mp\alpha/2}
    \end{bmatrix}
\end{equation*}

\subsubsection{Derivation: Spin operator in matrix representation.}
We start from the action of ladder operator
\begin{equation*}
    \hat{L}_{\pm}   |l, m\rangle = \hbar \sqrt{l(l+1) - m(m \pm 1)}   |l, m \pm 1\rangle
\end{equation*}
where $l=1/2$ and $m_s=\pm1/2$.
From the action of spin operator
\begin{equation*}
    S_+ |\downarrow\rangle = \frac{1 }{2 } |\uparrow\rangle, \quad S_- |\uparrow\rangle = \frac{1 }{2} |\downarrow\rangle
\end{equation*}
we write
\begin{align*}
    (S^+ + S^-) \begin{bmatrix}1 \\ 0 \end{bmatrix} & =  2 S_1 \begin{bmatrix}1 \\ 0 \end{bmatrix} = \begin{bmatrix}0 \\ 1 \end{bmatrix} \\
    (S^+ + S^-) \begin{bmatrix}0 \\ 1 \end{bmatrix} & =  2 S_1 \begin{bmatrix}0 \\ 1 \end{bmatrix} = \begin{bmatrix}1 \\ 0 \end{bmatrix}
\end{align*}
with $S_++S_-=S_1+iS_2-S_1-iS_2$.
Equivalently, writing the first equation as
\begin{equation*}
    2S_1 \begin{bmatrix}
        1 \\
        0
    \end{bmatrix}  =
    2
    \begin{bmatrix}
        a & b \\
        c & d
    \end{bmatrix}  \begin{bmatrix}
        1 \\
        0
    \end{bmatrix}  =
    2
    \begin{bmatrix}
        a \\
        c
    \end{bmatrix}
    =
    \begin{bmatrix}
        0 \\
        1
    \end{bmatrix}
\end{equation*}
we obtain $a=0$ and $c=1/2$.
Then writing the second equation in the same way
\begin{equation*}
    2S_1 \begin{bmatrix}
        0 \\
        1
    \end{bmatrix}  =
    \frac{1 }{2}
    \begin{bmatrix}
        a & b \\
        c & d
    \end{bmatrix}  \begin{bmatrix}
        0 \\
        1
    \end{bmatrix}  =
    2
    \begin{bmatrix}
        b \\
        d
    \end{bmatrix}
    =
    \begin{bmatrix}
        1 \\
        0
    \end{bmatrix}
\end{equation*}
we obtain $b=1/2$ and $d=0$.
The matrix representation is then
\begin{equation*}
    S_1 =
    \begin{bmatrix}
        a & b \\
        c & d
    \end{bmatrix}=
    \frac{1  }{2}
    \begin{bmatrix}
        0 & 1 \\
        1 & 0
    \end{bmatrix}
\end{equation*}

Next for the lowering operator
\begin{align*}
    (S^+ - S^-) \begin{bmatrix}1 \\ 0 \end{bmatrix} & =  2i S_2 \begin{bmatrix}1 \\ 0 \end{bmatrix} = - \begin{bmatrix}0 \\ 1 \end{bmatrix} \\
    (S^+ - S^-) \begin{bmatrix}0 \\ 1 \end{bmatrix} & =  2i S_2 \begin{bmatrix}0 \\ 1 \end{bmatrix} = \begin{bmatrix}1 \\ 0 \end{bmatrix}
\end{align*}
Equivalently, writing the first equation as
\begin{equation*}
    2iS_2 \begin{bmatrix}
        1 \\
        0
    \end{bmatrix}  =
    2i
    \begin{bmatrix}
        a & b \\
        c & d
    \end{bmatrix}  \begin{bmatrix}
        1 \\
        0
    \end{bmatrix}  =
    2i
    \begin{bmatrix}
        a \\
        c
    \end{bmatrix}
    =
    -
    \begin{bmatrix}
        0 \\
        1
    \end{bmatrix}
\end{equation*}
we obtain $a=0$ and $c=i/2$.
Then writing the second equation in the same way
\begin{equation*}
    2iS_1 \begin{bmatrix}
        0 \\
        1
    \end{bmatrix}  =
    2i
    \begin{bmatrix}
        a & b \\
        c & d
    \end{bmatrix}  \begin{bmatrix}
        0 \\
        1
    \end{bmatrix}  =
    2i
    \begin{bmatrix}
        b \\
        d
    \end{bmatrix}
    =
    \begin{bmatrix}
        1 \\
        0
    \end{bmatrix}
\end{equation*}
we obtain $b=-i/2$ and $d=0$.
The matrix representation is then
\begin{equation*}
    S_2 =
    \begin{bmatrix}
        a & b \\
        c & d
    \end{bmatrix}=
    \frac{1  }{2}
    \begin{bmatrix}
        0 & -i \\
        i & 0
    \end{bmatrix}
\end{equation*}

Finally, the matrix for $S_3$ produces the eigenvalues $m_s = \pm \frac{1}{2}$:
\begin{equation*}
    S_3 \left| \frac{1}{2}, \pm \frac{1}{2} \right\rangle = \pm \frac{1}{2} \left| \frac{1}{2}, \pm \frac{1}{2} \right\rangle
\end{equation*}
As such, we have
\begin{equation*}
    S_3 = \frac{1}{2} \begin{bmatrix}
        1 & 0  \\
        0 & -1
    \end{bmatrix}
\end{equation*}

\subsubsection{Derivation: Ladder operator matrix representation.}
For the increasing operator
\begin{equation*}
    \sigma_+
    = \frac{1}{2} \left( \begin{bmatrix} 0 & 1 \\ 1 & 0 \end{bmatrix} + i \begin{bmatrix} 0 & -i \\ i & 0 \end{bmatrix} \right)
    = \frac{1}{2} \left( \begin{bmatrix} 0 & 1 \\ 1 & 0 \end{bmatrix} + \begin{bmatrix} 0 & 1 \\ -1 & 0 \end{bmatrix}\right)
    = \begin{bmatrix} 0 & 1 \\ 0 & 0 \end{bmatrix}
\end{equation*}
and for the lowering
\begin{equation*}
    \sigma_-
    = \frac{1}{2} \left( \begin{bmatrix} 0 & 1 \\ 1 & 0 \end{bmatrix} - i \begin{bmatrix} 0 & -i \\ i & 0 \end{bmatrix} \right)
    = \frac{1}{2} \left( \begin{bmatrix} 0 & 1 \\ 1 & 0 \end{bmatrix} - \begin{bmatrix} 0 & 1 \\ -1 & 0 \end{bmatrix}\right)
    = \begin{bmatrix} 0 & 0 \\ 1 & 0 \end{bmatrix}
\end{equation*}

\subsubsection{Derivation: Pauli matrices identity.}
For the first case, we can prove them individually
\begin{align*}
    \sigma_1^2 & =
    \begin{bmatrix}
        0 & 1 \\
        1 & 0
    \end{bmatrix}
    \begin{bmatrix}
        0 & 1 \\
        1 & 0
    \end{bmatrix}
    =
    \begin{bmatrix}
        1 & 0 \\
        0 & 1
    \end{bmatrix}
    = \mathbf{1}   \\
    \sigma_2^2 & =
    \begin{bmatrix}
        0 & -i \\
        i & 0
    \end{bmatrix}
    \begin{bmatrix}
        0 & -i \\
        i & 0
    \end{bmatrix}
    =
    \begin{bmatrix}
        1 & 0 \\
        0 & 1
    \end{bmatrix}
    = \mathbf{1}   \\
    \sigma_3^2 & =
    \begin{bmatrix}
        1 & 0  \\
        0 & -1
    \end{bmatrix}
    \begin{bmatrix}
        1 & 0  \\
        0 & -1
    \end{bmatrix}
    =
    \begin{bmatrix}
        1 & 0 \\
        0 & 1
    \end{bmatrix}
    = \mathbf{1}
\end{align*}
Next we consider the case
\begin{equation*}
    \sigma_1 \sigma_2 = \begin{bmatrix} 0 & 1 \\ 1 & 0 \end{bmatrix} \begin{bmatrix} 0 & -i \\ i & 0 \end{bmatrix} = \begin{bmatrix} i & 0 \\ 0 & -i \end{bmatrix} = i \sigma_3
\end{equation*}
Instead of considering other case of multiplucation, we can cyclically permutate the index and obtain
\begin{equation*}
    \sigma_2 \sigma_3 = i \sigma_1, \quad \sigma_3 \sigma_1 = i \sigma_2
\end{equation*}
or $\sigma_{i}\sigma_j=i\epsilon_{ijk}\sigma_j$.
By combining both case, we have
\begin{equation*}
    \sigma_i \sigma_j = \delta_{ij} \mathbf{1} + i \epsilon_{ijk} \sigma_k
\end{equation*}

Using this, we prove the commutator relation
\begin{align*}
    [\sigma_i,\sigma_j]
     & =   \sigma_i \sigma_j - \sigma_j \sigma_i                                                                    \\
     & = (\delta_{ij} \mathbf{1} + i \epsilon_{jk} \sigma_k) - (\delta_{ji} \mathbf{1} + i \epsilon_{jik} \sigma_k) \\
     & =i \epsilon_{jk} \sigma_k - i \epsilon_{jik} \sigma_k = i \epsilon_{jk} \sigma_k + i \epsilon_{ijk} \sigma_k \\
    [\sigma_i,\sigma_j]
     & = 2i \epsilon_{jk} \sigma_k
\end{align*}
and the anticommutator
\begin{align*}
    \left\{ \sigma_i,\sigma_j \right\}
     & =   \sigma_i \sigma_j + \sigma_j \sigma_i                                                                    \\
     & = (\delta_{ij} \mathbf{1} + i \epsilon_{jk} \sigma_k) + (\delta_{ji} \mathbf{1} + i \epsilon_{jik} \sigma_k) \\
     & =\delta_{ij} \mathbf{1} + \delta_{ji} \mathbf{1} = \delta_{ij} \mathbf{1} + \delta_{ij} \mathbf{1}           \\
    \left\{ \sigma_i,\sigma_j \right\}
     & = 2 \delta_{ij} \mathbf{1}
\end{align*}

For the third identity, note that we can write the left term as
\begin{align*}
    (\mathbf{a} \cdot \boldsymbol{\sigma})(\mathbf{b} \cdot \boldsymbol{\sigma})
     & =  a_i b_j \sigma_i \sigma_j =     a_i b_j\delta_{ij} \mathbf{1} + ia_i b_j \epsilon_{ijk} \sigma_k \\
     & =  a_i b_i \mathbf{1} + i (\mathbf{a }\times \mathbf{b })_k \sigma_k                                \\
    (\mathbf{a} \cdot \boldsymbol{\sigma})(\mathbf{b} \cdot \boldsymbol{\sigma})
     & = (\mathbf{a }\cdot \mathbf{b }) + i (\mathbf{a }\times \mathbf{b } )\cdot \boldsymbol{\sigma}
\end{align*}

\subsubsection{Derivation: commutator and anticommutator relation for Pauli matrices ladder operator.}
First the commutator for third Pauli matrices and the ladder operator
\begin{align*}
    [\sigma_3, \sigma_\pm]
     & =  \frac{1}{2} \left( [\sigma_3, \sigma_1] \pm i[\sigma_3, \sigma_2] \right)                                      \\
     & =  \frac{1}{2} \left( 2i\sigma_2 \pm i(-2i\sigma_1) \right) = \frac{1}{2} \left( 2i\sigma_2 \pm 2\sigma_1 \right) \\
    [\sigma_3, \sigma_\pm]
     & = (\pm \sigma_1 \pm i \sigma_2) = \pm \sigma_{\pm}
\end{align*}
then the ladder operator with each other
\begin{align*}
    [\sigma_+,\sigma_-]
     & =  \frac{1}{4} [(\sigma_1 + i\sigma_2), (\sigma_1 - i\sigma_2)]                                                       \\
     & =  \frac{1}{4} \left( [\sigma_1,\sigma_1] - i[\sigma_1,\sigma_2] + i[\sigma_2,\sigma_1] + [\sigma_2,\sigma_2] \right) \\
     & =  \frac{1}{4}\left( -i(2i\sigma_3) + i(-2i\sigma_3) \right)                                                          \\
    [\sigma_+, \sigma_-]
     & = \frac{1}{4}(2\sigma_3 + 2\sigma_3) = \sigma_3
\end{align*}

For the anticommutator relation
\begin{align*}
    \{\sigma_{+}, \sigma_{-}\}
     & = \frac{1}{4}(\sigma_{x} + i \sigma_{y})(\sigma_{x} - i \sigma_{y}) + \frac{1}{4}(\sigma_{x} - i \sigma_{y})(\sigma_{x} + i \sigma_{y}) \\
     & =  \frac{1}{4} \left( \sigma_1 ^2- i \sigma_1 \sigma_2 + i \sigma_2 \sigma_1  + \sigma_2^2 \right)                                      \\
     & \qquad + \frac{1}{4} \left( \sigma_{x}^{2} + i \sigma_{x} \sigma_{y} - i \sigma_{y} \sigma_{x} + \sigma_{y}^{2} \right)                 \\
    \{\sigma_{+}, \sigma_{-}\}
     & = \frac{1 }{4 }\left( \mathbf{1}+\mathbf{1}+\mathbf{1}+\mathbf{1} \right) = \mathbf{1}
\end{align*}
and
\begin{align*}
    \{ \sigma_3, \sigma_\pm \}
     & =  \sigma_3 \sigma_\pm + \sigma_\pm \sigma_3                                                                        \\
     & = \frac{1}{2} (\sigma_3 \sigma_1 \pm i \sigma_3 \sigma_2) + \frac{1}{2} (\sigma_1 \sigma_3 \pm i \sigma_2 \sigma_3) \\
     & =  \frac{1}{2} (\sigma_3 \sigma_1 + \sigma_1 \sigma_3 \pm i (\sigma_3 \sigma_2 + \sigma_2 \sigma_3))                \\
     & = \frac{1 }{2 }\left( \left\{ \sigma_3,\sigma_1 \right\}\pm i \left\{ \sigma_3,\sigma_2 \right\}   \right)          \\
    \{ \sigma_3, \sigma_{\pm} \}
     & =   \frac{1}{2} (0 \pm i \cdot 0) = 0
\end{align*}


\subsubsection{Prove of how Pauli matrices span the $2 \times 2$ Hermitian matrices space.}
Start from arbitrary matrices
\begin{equation*}
    H = \begin{bmatrix}
        a & b \\
        c & d
    \end{bmatrix},
    \quad
    H ^\dagger = \begin{bmatrix}
        a ^* & c ^* \\
        b ^* & d ^*
    \end{bmatrix},
    \quad a, b, c, d \in \mathbb{C}
\end{equation*}
The Hermitian condition $H = H ^\dagger$ require
\begin{equation*}
    H = \begin{bmatrix}
        a   & b \\
        b^* & d
    \end{bmatrix}, \quad a, d \in \mathbb{R}, \quad b \in \mathbb{C}
\end{equation*}
The diagonal must be real since $a = a ^*$ and $d = d ^*$, meanwhile the diagonal can be complex with $c ^* = b$ implies $c  = b ^*$.
We assume that the space can be written as linear combination of
\begin{equation*}
    H =c_0 \mathbf{1 } + c_1 \sigma_1 + c_2 \sigma_2 + c_3 \sigma_3 =\begin{bmatrix}
        a   & b \\
        b^* & d
    \end{bmatrix}
\end{equation*}
such that
\begin{equation*}
    \begin{bmatrix}
        c_0 & 0   \\
        0   & c_0
    \end{bmatrix}
    +
    \begin{bmatrix}
        0   & c_1 \\
        c_1 & 0
    \end{bmatrix}
    +
    \begin{bmatrix}
        0    & -ic_2 \\
        ic_2 & 0
    \end{bmatrix}
    +
    \begin{bmatrix}
        c_3 & 0    \\
        0   & -c_3
    \end{bmatrix}
    =
    \begin{bmatrix}
        c_0 + c_3  & c_1 - ic_2 \\
        c_1 + ic_2 & c_0 - c_3
    \end{bmatrix}
\end{equation*}
Equating both, we obtain the results for diagonal elements
\begin{equation*}
    \begin{cases}
        c_0 + c_3 = a \\
        c_0 - c_3 = d
    \end{cases}
    \Rightarrow
    \quad
    \begin{cases}
        c_0 = (a+d)/2 \\
        c_3 = (a-d)/2
    \end{cases}
\end{equation*}
For the off diagonal, write $b=b_r+ib_i $
\begin{equation*}
    c_1 - i c_2 = b_r + i b_i
    \quad
    \Rightarrow
    \quad
    \begin{cases}
        c_1 = b_r \\
        c_2 = -b_i
    \end{cases}
\end{equation*}
All coefficients $c_0, c_1, c_2, c_3$ are uniquely determined, meaning there exists exactly one quadruple of the coefficient
\begin{equation*}
    c_0 = \frac{a + d}{2}, \quad c_3 = \frac{a - d}{2}, \quad c_1 = \mathrm{Re}(b), \quad c_2 = -\mathrm{Im}(b)
\end{equation*}
that reproduces the matrix $H$.
Thus, $\left\{ \mathbf{1 },\sigma_1 ,\sigma_2,\sigma_3\right\} $ span the space of all $2 \times 2$ Hermitian matrices.

\subsubsection{Derivation: Bloch vector representation of spinor state.}
For general normalized spinor state
\begin{equation*}
    |\chi\rangle = \begin{bmatrix} \alpha \\ \beta \end{bmatrix}, \quad |\alpha|^2 + |\beta|^2 = 1.
\end{equation*}
then the components of Bloch vector is
\begin{align*}
    e_x & =  \begin{bmatrix} \alpha & \beta \end{bmatrix}\begin{bmatrix} 0 & 1 \\ 1 & 0 \end{bmatrix} \begin{bmatrix} \alpha \\ \beta \end{bmatrix} = \alpha^*\beta + \beta^*\alpha = 2\text{Re}(\alpha^*\beta)      \\
    e_y & =  \begin{bmatrix} \alpha & \beta \end{bmatrix}\begin{bmatrix} 0 & -i \\ i & 0 \end{bmatrix} \begin{bmatrix} \alpha \\ \beta \end{bmatrix} = -i(\alpha^*\beta - \beta^*\alpha) = 2\text{Im}(\alpha^*\beta) \\
    e_z & =  \begin{bmatrix} \alpha & \beta \end{bmatrix} \begin{bmatrix} 1 & 0 \\ 0 & -1 \end{bmatrix}  \begin{bmatrix} \alpha \\ \beta \end{bmatrix} = |\alpha|^2 - |\beta|^2
\end{align*}
where we have used the complex relation $\mathrm{Re } (z) = (z + z ^* )/2$ and $\mathrm{Im } (z) = (z - z ^* )/2i$.
Computing the norm
\begin{align*}
    |\mathbf{e}|^2 & =  e_x^2 + e_y^2 + e_z^2 = 4\left[(\mathrm{Re}(\alpha^* \beta))^2 + (\mathrm{Im}(\alpha^* \beta))^2\right] + (|\alpha|^2 - |\beta|^2)^2 \\
                   & =  4|\alpha|^2 |\beta|^2 + |\alpha|^4 - 2|\alpha|^2|\beta|^2 + |\beta|^4 = |\alpha|^4 + 2 |\alpha|^2 |\beta|^2 + |\beta|^4              \\
    |\mathbf{e}|^2
                   & = (|\alpha|^2 + |\beta|^2)^2 = 1
\end{align*}
Hence the Bloch vector also normalized

Suppose we have the spinor state of
\begin{equation*}
    |\zeta\rangle = \frac{1}{\sqrt{2}} \begin{bmatrix}
        1 \\
        1
    \end{bmatrix}
\end{equation*}
All components of Bloch vector are
\begin{align*}
    e_x & =
    \begin{bmatrix}
        \frac{1}{\sqrt{2}} & \frac{1}{\sqrt{2}}
    \end{bmatrix}
    \begin{bmatrix} 0 & 1 \\ 1 & 0 \end{bmatrix}
    \begin{bmatrix} \frac{1}{\sqrt{2}} \\ \frac{1}{\sqrt{2}} \end{bmatrix} =
    \begin{bmatrix}
        \frac{1}{\sqrt{2}} & \frac{1}{\sqrt{2}}
    \end{bmatrix}
    \begin{bmatrix}
        \frac{1}{\sqrt{2}} \\ \frac{1}{\sqrt{2}}
    \end{bmatrix}
    = 1     \\
    e_y & =
    \begin{bmatrix}
        \frac{1}{\sqrt{2}} & \frac{1}{\sqrt{2}}
    \end{bmatrix}
    \begin{bmatrix} 0 & -i \\ i & 0 \end{bmatrix}
    \begin{bmatrix} \frac{1}{\sqrt{2}} \\ \frac{1}{\sqrt{2}} \end{bmatrix} =
    \begin{bmatrix}
        \frac{1}{\sqrt{2}} & \frac{1}{\sqrt{2}}
    \end{bmatrix}
    \begin{bmatrix}
        -\frac{i}{\sqrt{2}} \\ \frac{i}{\sqrt{2}}
    \end{bmatrix}
    = 0     \\
    e_z & =
    \begin{bmatrix}
        \frac{1}{\sqrt{2}} & \frac{1}{\sqrt{2}}
    \end{bmatrix}
    \begin{bmatrix} 1 & 0 \\ -0 & 1 \end{bmatrix}
    \begin{bmatrix} \frac{1}{\sqrt{2}} \\ \frac{1}{\sqrt{2}} \end{bmatrix} =
    \begin{bmatrix}
        \frac{1}{\sqrt{2}} & \frac{1}{\sqrt{2}}
    \end{bmatrix}
    \begin{bmatrix}
        \frac{1}{\sqrt{2}} \\ -\frac{1}{\sqrt{2}}
    \end{bmatrix}
    = 0
\end{align*}
So the Bloch vector $\mathbf{e}=(1,0,0)$ points along the positive $x$-axis on the Bloch sphere, where the spherical coordinates are
\begin{equation*}
    \theta = \arccos(e_z) = \frac{\pi}{2}\qquad
    \phi = \arctan2(e_y, e_x) = 0.
\end{equation*}
The spinor can be reconstructed from this
\begin{equation*}
    |\chi\rangle = \begin{bmatrix}
        \cos(\theta/2) \\
        e^{i\phi} \sin(\theta/2)
    \end{bmatrix}
    = \frac{1}{\sqrt{2}} \begin{bmatrix}
        1 \\
        1
    \end{bmatrix}
\end{equation*}

We can do the same thing with the computational base
\begin{equation*}
    \ket{0} = \begin{bmatrix}
        1 \\0
    \end{bmatrix}
    \qquad
    \ket{1} = \begin{bmatrix}
        0 \\1
    \end{bmatrix}
\end{equation*}
For the base $\ket{0}$
\begin{align*}
    e_x & =
    \begin{bmatrix} 1 & 0 \end{bmatrix}
    \begin{bmatrix} 0 & 1 \ 1 & 0 \end{bmatrix}
    \begin{bmatrix} 1 \\ 0 \end{bmatrix} =
    \begin{bmatrix} 1 & 0 \end{bmatrix}
    \begin{bmatrix} 0 \\ 1 \end{bmatrix} = 0 \\
    e_y & =
    \begin{bmatrix} 1 & 0 \end{bmatrix}
    \begin{bmatrix} 0 & -i \ i & 0 \end{bmatrix}
    \begin{bmatrix} 1 \\ 0 \end{bmatrix} =
    \begin{bmatrix} 1 & 0 \end{bmatrix}
    \begin{bmatrix} 0 \\ i \end{bmatrix} = 0 \\
    e_z & =
    \begin{bmatrix} 1 & 0 \end{bmatrix}
    \begin{bmatrix} 1 & 0 \\ 0 & -1 \end{bmatrix}
    \begin{bmatrix} 1 \\ 0 \end{bmatrix} =
    \begin{bmatrix} 1 & 0 \end{bmatrix}
    \begin{bmatrix} 1 \\ 0 \end{bmatrix} = 1
\end{align*}
and for $\ket{1}$
\begin{align*}
    e_x & =
    \begin{bmatrix} 0 & 1 \end{bmatrix}
    \begin{bmatrix} 0 & 1 \\ 1 & 0 \end{bmatrix}
    \begin{bmatrix} 0 \\ 1 \end{bmatrix} =
    \begin{bmatrix} 0 & 1 \end{bmatrix}
    \begin{bmatrix} 1 \\ 0 \end{bmatrix} = 0  \\
    e_y & =
    \begin{bmatrix} 0 & 1 \end{bmatrix}
    \begin{bmatrix} 0 & -i \\ i & 0 \end{bmatrix}
    \begin{bmatrix} 0 \\ 1 \end{bmatrix} =
    \begin{bmatrix} 0 & 1 \end{bmatrix}
    \begin{bmatrix} -i \\ 0 \end{bmatrix} = 0 \\
    e_z & =
    \begin{bmatrix} 0 & 1 \end{bmatrix}
    \begin{bmatrix} 1 & 0 \ 0 & -1 \end{bmatrix}
    \begin{bmatrix} 0 \\ 1 \end{bmatrix} =
    \begin{bmatrix} 0 & 1 \end{bmatrix}
    \begin{bmatrix} 0 \\ -1 \end{bmatrix} = -1
\end{align*}
As expected, $\ket{0}$ has the vector $\mathbf{e} = (0,0,1)$, while $\ket{1}$ has $\mathbf{e}=(0,0,-1)$.
The spherical coordinates for $\ket{0}$ are $(\theta,\phi)=(0,0)$, while for $\ket{1}$ are $(\theta,\phi)=(\pi,0)$.
The reconstructed basis takes the form
\begin{equation*}
    \ket{0}= \begin{bmatrix}
        1 \\0
    \end{bmatrix}
    \qquad
    \ket{1} = \begin{bmatrix}
        0 \\1
    \end{bmatrix}
\end{equation*}

\subsubsection{Derivation: Pauli matrices as projector operator.}
The functional relation for projection operator first derived by assuming that the function of operator $f(P)$ can be expanded as power series
\begin{equation*}
    f(P) = \sum_{n=0}^{\infty} a_n P^n
\end{equation*}
Idempotent $P ^{2 }= P$ implies that $P ^{n }=P$ for $n \geq1$, so
\begin{equation*}
    f(P) = a_0 I + \sum_{n=1}^{\infty} a_n P^n = a_0 I + \left( \sum_{n=1}^{\infty} a_n  \right)  P
\end{equation*}
Suppose we evaluate the scalar function $f(x)$ at $x=0$ and $x=1$
\begin{equation*}
    f(x) = \sum_{n=0}^{\infty} a_n x^n =
    \begin{cases}
        f(0) = a_0 \\
        f(1) = \sum_{n=0}^{\infty} a_n = a_0 + \sum_{n=1}^{\infty} a_n
    \end{cases}
\end{equation*}
if we define $0 ^{0}$.
Then we have the series as
\begin{equation*}
    \sum_{n=1}^{\infty} a_n = f(1) - f(0)
\end{equation*}
Substituting into the operator expression gives
\begin{equation*}
    f(P) = f(0)I + [f(1) - f(0)]P
\end{equation*}
Next step is simply use the projection operator $\sigma_+\sigma_-$ or $\sigma_-\sigma+$ to obtain the desired relation.

Next we derive the transformation by similarity for Pauli matrices.
First note that since $\sigma_3$ is diagonal, we can write the exponent of it by
\begin{equation*}
    \exp \left(\pm \frac{\alpha \sigma_z}{2}  \right)
    = \begin{bmatrix}
        e^{\pm\alpha/2} & 0               \\
        0               & e^{\mp\alpha/2}
    \end{bmatrix}
\end{equation*}
which I have proven somewhere.
In case you don't believe it, we can directly derive by writing the power series representation
\begin{equation*}
    \exp \left( \frac{\alpha \sigma_z}{2}  \right) = \sum_{n=0}^{\infty} \frac{1}{n!} \left( \frac{\alpha \sigma_3}{2} \right)^{n}
\end{equation*}
Since $\sigma_3^2 =I$, we have $\sigma_3^{2n} = I$ and $\sigma_3^{2n+1} = \sigma_3$.
As such, we split the series into even and odd power
\begin{equation*}
    \exp \left( \frac{\alpha \sigma_z}{2}  \right)
    = \sum_{n=0}^{\infty} \frac{1}{(2n)!} \left( \frac{\alpha}{2} \right)^{2n} I + \sum_{n=0}^{\infty} \frac{1}{(2n + 1)!} \left( \frac{\alpha}{2} \right)^{2n+1} \sigma_3
\end{equation*}
Then simply evaluate it
\begin{align*}
    \exp \left( \frac{\alpha \sigma_z}{2}  \right)
     & =  \cosh \left( \frac{\alpha }{2 }  \right) I + \sinh \left( \frac{\alpha }{2 }  \right) \sigma_3 \\
     & = \begin{bmatrix}
             \cosh(\alpha/2) + \sinh(\alpha/2) & 0                                 \\
             0                                 & \cosh(\alpha/2) - \sinh(\alpha/2)
         \end{bmatrix}                           \\
    \exp \left( \frac{\alpha \sigma_z}{2}  \right)
     & = \begin{bmatrix}
             e^{\alpha/2} & 0             \\
             0            & e^{-\alpha/2}
         \end{bmatrix}
\end{align*}
same goes for the negative exponent.
All that left is direct multiplication
\begin{align*}
    \exp \left( \frac{\alpha \sigma_z}{2} \right)  \sigma_+ \exp \left(- \frac{\alpha \sigma_z}{2} \right)
     & =
    \begin{bmatrix}
        e^{\alpha/2} & 0             \\
        0            & e^{-\alpha/2}
    \end{bmatrix}
    \begin{bmatrix}
        0 & 1 \\
        0 & 0
    \end{bmatrix}
    \begin{bmatrix}
        e^{-\alpha/2} & 0            \\
        0             & e^{\alpha/2}
    \end{bmatrix} \\
     & =
    \begin{bmatrix}
        e^{\alpha/2} & 0             \\
        0            & e^{-\alpha/2}
    \end{bmatrix}
    \begin{bmatrix}
        0 & e^{\alpha/2} \\
        0 & 0
    \end{bmatrix}             \\
     & = \begin{bmatrix}
             0 & e^{\alpha} \\
             0 & 0
         \end{bmatrix}
    = e^{\alpha} \sigma_+
\end{align*}
Once again, the same goes for the negative exponent.


\subsection{Dirac Equation}
The Dirac equation is
\begin{equation*}
    i \hbar \; \partial_t \Psi(\mathbf{r },t) = H \Psi(\mathbf{r },t)
\end{equation*}
with the Hamiltionian
\begin{equation*}
    H = c \hat{\boldsymbol{\alpha}} \cdot \mathbf{p} + \hat{\beta}mc^2
    = \begin{bmatrix}
        m c^{2} \mathbf{1}                        & c\,\boldsymbol{\sigma}\cdot\hat{\mathbf{p}} \\
        c\boldsymbol{\sigma}\cdot\hat{\mathbf{p}} & -m c^{2} \mathbf{1}
    \end{bmatrix}
\end{equation*}
and $\Psi(\mathbf{r },t)$ as the four-component spinor. 
It can also be written in this way 
\begin{equation*}
    i\hbar \;\partial_t \Psi(\mathbf{r},t)=\left(-i\hbar c\hat{\boldsymbol{\alpha}}\cdot\nabla + \hat{\beta}m c^{2}\right)\Psi(\mathbf{r},t)
\end{equation*}

The matrices $\hat{\alpha}_i$ and $\hat{\beta}$ defined as
\begin{equation*}
    {\hat{\alpha}}_i = \begin{bmatrix} \mathbf{0} & {\sigma}_i\\  {\sigma}_i & \mathbf{0} \end{bmatrix}\qquad
    \beta = \begin{bmatrix} \mathbf{1} & \mathbf{0}\\ \mathbf{0} & -\mathbf{1} \end{bmatrix}
\end{equation*}
and satisfies
\begin{equation*}
    \{\hat{\alpha}_i,\alpha_j\}=2\delta_{ij} \qquad
    \{\hat{\alpha}_i,\beta\}=\mathbf{0} \qquad
    \hat{\beta}^2=\mathbf{0}
\end{equation*}

The equation has the same structure with the Schrödinger equation, in fact this is the relativistic generalization of the Schrödinger equation.
This equation is constructed specifically to satisfy the relativistic energy $E= \sqrt{\mathbf{p }^2c^2+m^2c^4}$.

\subsubsection{Spinor-field.}
The four-components spinor encode two states of spin-$1/2$ particle and degenerate solution of relativistic dispersion relation $E^2 = p^2c^2+m^2c^4$.
The components can be grouped as 
\begin{equation*}
\Psi = \begin{bmatrix}\phi\\\chi\end{bmatrix} \qquad
\phi = \begin{bmatrix}\psi_1\\\psi_2\end{bmatrix} \qquad
\chi = \begin{bmatrix}\psi_3\\\psi_4\end{bmatrix}
\end{equation*}
$\phi$ component encode degrees of freedom at low energy; while $\chi$ encode relativistic correction responsible for spin–orbit coupling, magnetic moment, relativistic effects.
The coupling of $\phi$ and $\chi$
\begin{equation*}
 \begin{bmatrix}
        m c^{2}                         & c\,\boldsymbol{\sigma}\cdot\hat{\mathbf{p}} \\
        c\boldsymbol{\sigma}\cdot\hat{\mathbf{p}} & -m c^{2} 
    \end{bmatrix}
\begin{bmatrix}\phi\\\chi\end{bmatrix}
= \begin{bmatrix}
mc^{2}\phi + c\bigl(\boldsymbol{\sigma}\cdot\mathbf{p}\bigr)\chi\\
c\bigl(\boldsymbol{\sigma}\cdot\mathbf{p}\bigr)\phi - mc^{2}
\end{bmatrix}
\end{equation*}
via $\boldsymbol{\sigma}\cdot \mathbf{p}$  is what makes the theory relatGivistic.

\subsubsection{Klein-Gordon equation.}
The Klein-Gordon originaly interpreted to generalize the Schrödinger equation in relativistic.
It starts with the relativistic energy
\begin{equation*}
    E^2 = p^2c^2+ m^2c^4
\end{equation*}
then introduce the prescription $    \mathbf{p} \to -i\hbar\nabla \qquad$ and  $    H \to i\hbar \;\partial_t$
\begin{equation*}
    \bigg(\frac{1}{c^2}\frac{\partial^2}{\partial t^2} - \nabla^2 + \frac{m^2 c^2}{\hbar^2}\bigg)\psi(\mathbf r,t) = 0
\end{equation*}
The equation is also known as the generalization of Maxwell’s wave equation for massive particles.
Asidem, the equation suffer two problem.
\begin{enumerate}
    \item \textbf{Negative energy.}
    The equation was supposed to describe single relativistic particle $\psi(\mathbf{r},t)$, it however include both spectrum of $+E$ and $-E$, so the theory is unbounded below if interpreted as a single-particle system. 

    In relativistic quantum field theory, this is corrected by interpreting $\psi(\mathbf{r},t)$ as the quantum field, and the negative energy as antiparticle. 
    \item \textbf{Incorrect probability density.} 
    A probability density must satisfies conserved
    \begin{equation*}
        \frac{d }{dt } \int d^3 \mathbf{r} \rho(\mathbf{r},t) = 0
    \end{equation*} 
    and positive definite
    \begin{equation*}
        \rho(\mathbf{r},t) \geq  0
    \end{equation*}
    quantity.
    Single plane wave state $\rho(\mathbf{r},t) = |\psi|^2$ does produce time-independent $\rho(\mathbf{r})$, however superpositions of pure states cause interference term to appear and time dependence to be produced.
    There exists a conserved current
    \begin{equation*}
        \rho = \frac{i\hbar}{2 m c^{2}}\bigl(\psi^{*}\partial_{t}\psi - \psi\partial_{t}\psi^{*}\bigr)
    \end{equation*}
    that satisfies the time-independence.
    This probability density, however, is not positive definite, so it still invalidates its interpretation as a probability density.
\end{enumerate}

Dirac's equation able to produce the continuity equation
\begin{equation*}
    \partial_t \rho(\mathbf{r },t) + c \nabla \cdot \mathbf{j}(\mathbf{r },t)=0
\end{equation*}
with $\rho(\mathbf{r },t)$ and $\mathbf{j}(\mathbf{r },t)$ three-valued 
\begin{equation*}
    \rho  (\mathbf{r },t)= \Psi ^\dagger(\mathbf{r },t)\Psi(\mathbf{r },t) \qquad
    \mathbf{j}(\mathbf{r },t)=\Psi ^\dagger (\mathbf{r },t)\hat{\boldsymbol{\alpha } }\Psi(\mathbf{r },t)
\end{equation*}
which does not suffer from non-conserved probability density
\begin{equation*}
\partial_t \int d^3r\;\rho =0
\end{equation*}

\subsubsection{Derivation: Dirac approach.}
He considered a wave equation with only first derivatives with respect to time and space and introduced the classical energy
\begin{equation*}
    E = c\boldsymbol{\hat{\alpha}}\cdot\mathbf{p} + \hat{\beta}m c^2
\end{equation*}
such that squaring it one recovers the relativistic energy
\begin{equation*}
    E ^{2 } =  \mathbf{p }^2c^2 +m^2c^4
\end{equation*}
Since spatial translations are generated by momentum and time translations are generated by the Hamiltonian, we introduce quantization prescription
\begin{equation*}
    \mathbf{p} \to -i\hbar\nabla \qquad
    H \to i\hbar \;\partial_t
\end{equation*}
The classical energy now reads
\begin{equation*}
    i\hbar \;\partial_t \Psi(\mathbf{r},t)=\left(-i\hbar c\hat{\boldsymbol{\alpha}}\cdot\nabla + \hat{\beta}m c^{2}\right)\Psi(\mathbf{r},t)
\end{equation*}
Notice that the wavefunction $\Psi(\mathbf{r },t)$, has four components in the abstract space of Dirac matrices, as such
\begin{equation*}
    i\hbar\partial_t
    \begin{bmatrix}
        \psi_1(\mathbf{r},t) \\
        \psi_2(\mathbf{r},t) \\
        \psi_3(\mathbf{r},t) \\
        \psi_4(\mathbf{r},t)
    \end{bmatrix}
    =
    \hat{H}
    \begin{bmatrix}
        \psi_1(\mathbf{r},t) \\
        \psi_2(\mathbf{r},t) \\
        \psi_3(\mathbf{r},t) \\
        \psi_4(\mathbf{r},t)
    \end{bmatrix}
\end{equation*}

The first term of the Hamiltonian has the elements of 
\begin{align*}
    \boldsymbol{\sigma}\cdot \hat{\mathbf p}
     & = -i\hbar\left( \sigma_x \partial_x + \sigma_y \partial_y + \sigma_z \partial_z \right) \\
     & = -i\hbar \left(
    \begin{bmatrix}
            0          & \partial_x \\
            \partial_x & 0
        \end{bmatrix}
    +
    \begin{bmatrix}
            0           & -i\partial_y \\
            i\partial_y & 0
        \end{bmatrix}
    +
    \begin{bmatrix}
            \partial_z & 0 \\
            0          & -\partial_z
        \end{bmatrix}
    \right) \\
     & = -i\hbar
    \begin{bmatrix}
        \partial_z               & \partial_x - i\partial_y \\
        \partial_x + i\partial_y & -\partial_z
    \end{bmatrix}
\end{align*}
Along the second term being trivial, the Hamiltonian explicitly written as 
\begin{equation*}
    \begin{bmatrix}
        mc^2                                             & 0                                                & -i\hbar c\;\partial_z                            & -i \hbar c \left( \partial_x-i \;\partial_y \right) \\
        0                                                & mc^2                                             & -i\hbar c\left(\partial_x + i\;\partial_y\right) & i\hbar c\,\partial_z                                \\
        -i\hbar c\;\partial_z                            & -i\hbar c\left(\partial_x - i\;\partial_y\right) & -mc^2                                            & 0                                                   \\
        -i\hbar c\left(\partial_x + i\;\partial_y\right) & i\hbar c\;\partial_z                             & 0                                                & -mc^2
    \end{bmatrix}
\end{equation*}

\subsubsection{Derivation: continuity equation.}
Start from the Dirac equation and its conjugate
\begin{gather*}
    \partial_t \Psi(\mathbf{r},t)=-c\hat{\boldsymbol{\alpha}}\cdot\nabla \Psi(\mathbf{r},t) 
     + (i m c^{2}/\hbar )\hat{\beta} \Psi(\mathbf{r},t) \\
    \partial_t \Psi ^\dagger(\mathbf{r},t)=-c [\nabla \Psi ^\dagger(\mathbf{r},t)] \cdot \hat{\boldsymbol{\alpha}}
    -(i m c^{2}/\hbar )\Psi ^\dagger(\mathbf{r},t) \hat{\beta} 
\end{gather*}
Left multiply the Dirac equation with the Hermitian conjugate of the wavefunction
\begin{equation*}
    \Psi ^\dagger \partial_t \Psi=c\Psi ^\dagger\hat{\boldsymbol{\alpha}}\cdot\nabla \Psi 
    + (i m c^{2}/\hbar ) \Psi ^\dagger \hat{\beta} \Psi 
\end{equation*}
the right multiply the Hermitian conjugate of the Dirac equation with the wavefunction
\begin{equation*}
    (\partial_t \Psi ^\dagger) \Psi=c [\nabla \Psi ^\dagger] \cdot \hat{\boldsymbol{\alpha}}\Psi
    -(i m c^{2}/\hbar )\Psi ^\dagger \hat{\beta} \Psi
\end{equation*}
and add both 
\begin{equation*}
    \partial_t \Psi ^\dagger\Psi 
    = -c \left[     \Psi ^\dagger\hat{\boldsymbol{\alpha}}\cdot\nabla \Psi + (\nabla \Psi ^\dagger) \cdot \hat{\boldsymbol{\alpha}}\Psi \right] 
\end{equation*}

Noting that $\hat{\boldsymbol{\alpha}}$ is constant in space such that $\nabla \hat{\boldsymbol{\alpha}} = 0$, can write
\begin{align*}
    \nabla \cdot (\Psi ^\dagger \hat{\boldsymbol{\alpha } }\Psi) 
    & = \sum_{i } \partial_i\big(\Psi^\dagger\hat{\alpha}_i\Psi\big)\\
    & =  \sum_{i } (\partial_i\Psi^\dagger)\,\hat{\alpha}_i\,\Psi + \Psi^\dagger(\partial_i\hat{\alpha}_i)\Psi + \Psi^\dagger\hat{\alpha}_i(\partial_i\Psi)\\
    & = \sum_i\big[(\partial_i\Psi^\dagger)\hat{\alpha}_i\,\Psi + \Psi^\dagger\hat{\alpha}_i(\partial_i\Psi)\big]\\
    \nabla \cdot (\Psi ^\dagger \hat{\boldsymbol{\alpha } }\Psi) 
    & = (\nabla\Psi^\dagger)\cdot\boldsymbol{\hat{\alpha}}\Psi + \Psi^\dagger\big(\boldsymbol{\hat{\alpha}}\cdot\nabla\Psi\big)
\end{align*}
So we obtain the continuity equation
\begin{equation*}
    \partial_t \rho + c \nabla \cdot \mathbf{j}=0
\end{equation*}
where 
\begin{equation*}
    \rho  = \Psi ^\dagger\Psi \qquad
    \mathbf{j}=\Psi ^\dagger \hat{\boldsymbol{\alpha } }\Psi
\end{equation*}

Now integrate over all space
\begin{equation*}
\int d^3r\;\partial_t \rho + \int d^3r\;\nabla \cdot  \mathbf{j} = 0    
\end{equation*}
Using divergence theorem 
\begin{equation*}
\partial_t \int d^3r\;\rho + \oint d \mathbf{a}\cdot  \mathbf{j} = 0        
\end{equation*}
Applying vanishing current at infinity, we obtain the conserved probability density
\begin{equation*}
\partial_t \int d^3r\;\rho =0
\end{equation*}
\end{document}